\begin{thebibliography}{99}

\bibitem{topham2015}
Topham C, Tighe A, Ly P, et al. MYC Is a Major Determinant of Mitotic Cell
Fate. \emph{Cancer Cell}. 2015;28:129--140.
\url{https://doi.org/10.1016/j.ccell.2015.06.001}.

\bibitem{pillay2019}
Pillay N, Tighe A, Nelson L, et al. DNA Replication Vulnerabilities Render
Ovarian Cancer Cells Sensitive to Poly(ADP-Ribose) Glycohydrolase Inhibitors.
\emph{Cancer Cell}. 2019;35:519--533.e8.
\url{https://doi.org/10.1016/j.ccell.2019.02.004}.

\bibitem{nelson2020}
Nelson L, Tighe A, Golder A, et al. A living biobank of ovarian cancer ex vivo
models reveals profound mitotic heterogeneity. \emph{Nature Communications}.
2020;11:822. \url{https://doi.org/10.1038/s41467-020-14551-2}.

\bibitem{williams2020}
Williams MS, Basma NJ, Amaral FMR, et al. Targeted nanopore sequencing for the
identification of ABCB1 promoter translocations in cancer. \emph{BMC Cancer}.
2020;20:1075. \url{https://doi.org/10.1186/s12885-020-07571-0}.

\bibitem{bronder2021}
Bronder D, Tighe A, Wangsa D, et al. TP53 loss initiates chromosomal
instability in fallopian tube epithelial cells. \emph{Disease Models \&
Mechanisms}. 2021;14:dmm049001.
\url{https://doi.org/10.1242/dmm.049001}.

\bibitem{barnes2021}
Barnes BM, Nelson L, Tighe A, et al. Distinct transcriptional programs
stratify ovarian cancer cell lines into the five major histological subtypes.
\emph{Genome Medicine}. 2021;13:140.
\url{https://doi.org/10.1186/s13073-021-00952-5}.

\bibitem{coulsongilmer2021}
Coulson-Gilmer C, Morgan RD, Nelson L, et al. Replication catastrophe is
responsible for intrinsic PAR glycohydrolase inhibitor-sensitivity in
patient-derived ovarian cancer models. \emph{Journal of Experimental \&
Clinical Cancer Research}. 2021;40:323.
\url{https://doi.org/10.1186/s13046-021-02124-0}.

\bibitem{golder2022}
Golder A, Nelson L, Tighe A, et al. Multiple-low-dose therapy: effective
killing of high-grade serous ovarian cancer cells with ATR and CHK1
inhibitors. \emph{NAR Cancer}. 2022;4:zcac036.
\url{https://doi.org/10.1093/narcan/zcac036}.

\bibitem{nelson2023}
Nelson L, Barnes BM, Tighe A, et al. Exploiting a living biobank to delineate
mechanisms underlying disease-specific chromosome instability.
\emph{Chromosome Research}. 2023;31:21.
\url{https://doi.org/10.1007/s10577-023-09731-x}.

\bibitem{morgan2023}
Morgan RD, Clamp AR, Barnes BM, et al. Homologous recombination deficiency in
newly diagnosed FIGO stage III/IV high-grade epithelial ovarian cancer: a
multi-national observational study. \emph{International Journal of
Gynecological Cancer}. 2023;33:1253--1259.
\url{https://doi.org/10.1136/ijgc-2022-004211}.

\bibitem{meeson2024}
Meeson KE, Schwartz J-M. Constraint-based modelling predicts metabolic
signatures of low- and high-grade serous ovarian cancer. \emph{npj Systems
Biology and Applications}. 2024;10:96.
\url{https://doi.org/10.1038/s41540-024-00418-5}.

\bibitem{coulsongilmer2024}
Coulson-Gilmer C, Littler S, Barnes BM, et al. Intrinsic PARG inhibitor
sensitivity is mimicked by TIMELESS haploinsufficiency and rescued by
nucleoside supplementation.
\emph{NAR Cancer}. 2024;6:zcae030.
\url{https://doi.org/10.1093/narcan/zcae030}.

\bibitem{littler2025}
Littler S, Barnes BM, Owen R, et al. Targeting SUMOylation in ovarian cancer:
sensitivity, resistance, and the role of MYC. \emph{iScience}.
2025;28:112555. \url{https://doi.org/10.1016/j.isci.2025.112555}.

\bibitem{tighe2025}
Tighe A, Nelson L, Morgan RD, et al. Screening a living biobank identifies
cabazitaxel as a strategy to combat acquired taxol resistance in high-grade
serous ovarian cancer. \emph{Cell Reports Medicine}. 2025;6:102160.
\url{https://doi.org/10.1016/j.xcrm.2025.102160}.

\bibitem{morgan2025}
Morgan RD, Burghel GJ, Shaw J, et al. Homologous recombination deficiency in
unselected cases of high-grade ovarian carcinoma. \emph{Journal of Medical
Genetics}. 2025; advance online publication.
\url{https://doi.org/10.1136/jmg-2025-110903}.

\bibitem{meeson2026}
Meeson KE, McGrail JC, Schwartz J-M, Taylor SS. Transcriptome-driven
constraint-based modelling reveals metabolic targets for ovarian cancer.
\emph{Cancer \& Metabolism}. 2026;14:7.
\url{https://doi.org/10.1186/s40170-026-00425-6}.

\bibitem{liu2019}
Liu A, Trairatphisan P, Gjerga E, et al. From expression footprints to causal
pathways: contextualizing large signaling networks with CARNIVAL.
\emph{npj Systems Biology and Applications}. 2019;5:40.
\url{https://doi.org/10.1038/s41540-019-0109-6}.

\bibitem{corneto}
Saez-Rodriguez J, et al. CORNETO: Constraint-based Optimization for the
Reconstruction of NETworks from Omics. Documentation and software repository.
\url{https://corneto.org/latest/guide/signaling/carnival.html}.

\bibitem{omnipath}
Türei D, et al. Integrated intra- and intercellular signaling knowledge for
multicellular omics analysis. \emph{Molecular Systems Biology}. 2021;17:e9923.
\url{https://doi.org/10.15252/msb.20209923}.

\bibitem{collectri}
Müller-Dott S, Tsirvouli E, Vazquez M, et al. Expanding the coverage of
regulons from high-confidence prior knowledge for accurate estimation of
transcription factor activities. \emph{Nucleic Acids Research}.
2023;51:10934--10949. \url{https://doi.org/10.1093/nar/gkad841}.

\bibitem{human1}
Robinson JL, Kocabaş P, Wang H, et al. An atlas of human metabolism.
\emph{Science Signaling}. 2020;13(624):eaaz1482.
\url{https://doi.org/10.1126/scisignal.aaz1482}.

\bibitem{brunet2004}
Brunet J-P, Tamayo P, Golub TR, Mesirov JP. Metagenes and molecular pattern
discovery using matrix factorization. \emph{Proceedings of the National
Academy of Sciences USA}. 2004;101:4164--4169.
\url{https://doi.org/10.1073/pnas.0308531101}.

\bibitem{taylorprofile}
University of Manchester Research Explorer. Stephen Taylor: research outputs,
projects and ORCID. Accessed 12 August 2026.
\url{https://research.manchester.ac.uk/en/persons/stephen-taylor}.

\bibitem{golderthesis2021}
Golder A. \emph{Development of a drug profiling platform for patient-derived
ovarian tumour cultures}. PhD thesis, University of Manchester, 2021.
\url{https://research.manchester.ac.uk/en/studentTheses/development-of-a-drug-profiling-platform-for-patient-derived-ovar/}.

\bibitem{bronderthesis2021}
Bronder D. \emph{Modelling chromosomal instability in high-grade serous
ovarian cancer}. PhD thesis, University of Manchester, 2021.
\url{https://research.manchester.ac.uk/en/studentTheses/modelling-chromosomal-instability-in-high-grade-serous-ovarian-ca/}.

\bibitem{morganthesis2021}
Morgan RD. \emph{Novel therapeutic strategies to treat high-grade serous
ovarian carcinoma by targeting modulators of PAR dynamics}. PhD thesis,
University of Manchester, 2021.
\url{https://research.manchester.ac.uk/en/studentTheses/novel-therapeutic-strategies-to-treat-high-grade-serous-ovarian-c/}.

\bibitem{owenthesis2023}
Owen R. \emph{Exploiting SUMOylation to target MYC-driven cancers}. PhD thesis,
University of Manchester, 2023.
\url{https://research.manchester.ac.uk/en/studentTheses/exploiting-sumoylation-to-target-myc-driven-cancers}.

\bibitem{meesonthesis2024}
Meeson KE. \emph{Multi-omics modelling to identify actionable targets in
epithelial ovarian cancer}. PhD thesis, University of Manchester, 2024.
\url{https://research.manchester.ac.uk/en/studentTheses/multi-omics-modelling-to-identify-actionable-targets-in-epithelia/}.

\end{thebibliography}
